\documentclass[11pt]{article}
\usepackage[T1]{fontenc}
\usepackage{lmodern}
\usepackage[margin=1in]{geometry}
\usepackage{amsmath,amssymb,amsthm}
\usepackage{longtable,booktabs,array,calc}
\usepackage[expansion=false]{microtype}
\usepackage[hidelinks,hypertexnames=false]{hyperref}
\setlength{\emergencystretch}{2em}
\providecommand{\tightlist}{\setlength{\itemsep}{0pt}\setlength{\parskip}{0pt}}
\newcommand{\workbenchurl}{https://github.com/KokunoYumeto/zeta-and-l-functions-workbench}
\newcommand{\recordurl}{https://github.com/KokunoYumeto/mathematics-workbenches}
\title{Testing a quadratic-polynomial average\\against a balanced cubic convolution\\
\large An applicability analysis of Grimmelt--Merikoski--Pandey--Pascadi}
\author{}
\date{24 November 2026}
\begin{document}
\maketitle
\begin{abstract}
The quadratic-polynomial Möbius estimate of Grimmelt, Merikoski,
Pandey, and Pascadi (2026) is the starting point of this analysis.
We test its actual shift-weight uniformity and Type II range against
the balanced two-ideal convolution over the Eisenstein field.
Selecting square shifts incurs an explicit density loss.
Multiplication preserves a square shift, and the part falling
in the source's small-shift range admits an elementary
three-quarter-power counting bound at the product scale.
Differencing generates many shifts but introduces a two-Möbius
correlation with a varying quadratic base. These are quantified
failures of the specified applications, not an impossibility
theorem for new dispersion methods. No bound for the complete
retained cubic sum or priority claim is made.
The mathematics and text were produced by GPT-6 Astra in ChatGPT Pro
mode; no human has checked this text.
\end{abstract}

\section{Statements and scope}
We test the quadratic-polynomial estimates of Grimmelt, Merikoski,
Pandey, and Pascadi (2026), abbreviated GMPP, against a fixed
two-ideal cubic convolution. The proof below credits their theorem
and records its exact numerical range. It establishes failures of
specified applications and elementary bounds for specified regions.
It does not establish the impossibility of a redesigned dispersion
argument or a bound for the complete remaining signed sum.

Let \(K=\mathbf Q(\omega)\), \(\omega^2+\omega+1=0\), and let
\(\psi\) be either faithful cubic character attached to
\(K(\sqrt[3]{269})/K\). Delete prime ideals over 3 and 269.
Put \(m(\mathfrak b)=\mu_K(\mathfrak b)\psi(\mathfrak b)\).
An integral ideal is primitive if it is divisible by no rational
principal ideal \((r)\), \(r>1\).
Fix positive constants defining dyadic norm intervals of size \(Z\),
a bounded smooth product weight \(F\), and \(\mathcal X\asymp Z^2\).
The expression examined is
\[
 \mathcal B=\sum_{\substack{\mathfrak b_1,\mathfrak b_2\ {\rm primitive}\\
                    N\mathfrak b_1\asymp Z,\ N\mathfrak b_2\asymp Z}}
 m(\mathfrak b_1)m(\mathfrak b_2)
 F(N\mathfrak b_1N\mathfrak b_2/\mathcal X).
\]
The general balanced convolution retains this primitive corner.
No assertion that its other content blocks have balanced primitive
norms is intended.

\paragraph{Arithmetic and counting assertions.}
Every primitive ideal with nonzero Möbius coefficient, except the
unit ideal, has a unique generator \(x+y\sqrt{-3}\) with \(y>0\)
after choosing the unit class that is 1 modulo 2. Its norm is
\(x^2+3y^2\), with \((x,y)=1\) and \(x+y\) odd.
For \(H\ge1\), the number of pairs in the preceding balanced block
with \(3y_1^2\le H\) is \(O(Z^{3/2}(1+\sqrt H))\).
In particular, with \(H=Z^{1/2-\nu}\), it is
\(O(\mathcal X^{7/8-\nu/4}+\mathcal X^{3/4})\).
If both coordinate shifts satisfy this restriction, the bound is
\(O(\mathcal X^{3/4-\nu/2}+\mathcal X^{1/2})\).
These are absolute counting estimates for the indicated pieces.

Multiplication gives
\[
 (x_1^2+3y_1^2)(x_2^2+3y_2^2)
 =(x_1x_2+3y_1y_2)^2+3(x_1y_2-x_2y_1)^2.
\]
For \(w=x_1y_2-x_2y_1\), the number of primitive pairs with
\(3w^2\le H\) is \(O(Z(1+\sqrt H))\).
At \(H=\mathcal X^{1/2-\nu}\), this becomes
\(O(\mathcal X^{1/2}+\mathcal X^{3/4-\nu/2})\).
The constants depend only on the fixed dyadic intervals.

\paragraph{Sparse-shift and scale assertions.}
For an array \(|S_h|\le A\) satisfying
\(\sum_{h\le H}|S_h|\le EHA\), no cancellation on a prescribed
set of \(B\le EH\) shifts follows: the array
\(S_h=A\mathbf1_{\mathcal S}(h)\) is a counterexample to that
implication. Two copies of the same first-moment estimate do
not change this conclusion. For square shifts \(B\asymp H^{1/2}\),
GMPP's bound with ambient \(H=V^{1/2-\nu}\) loses the factor
\(\nu V^{1/4-\nu/2}\) relative to the selected support.
Its Type II saving \(V^{-\eta/400+o(1)}\) is smaller than this
density loss throughout its stated range
\(0\le\nu\le1/6-2\eta\).

Even if coefficient matching were granted, the product-norm
choice \(V=\mathcal X\), \(M\asymp N\asymp V^{1/2}\) is outside
GMPP's Type II condition
\(N<V^{1/4-\nu/2-\eta}=H^{1/2}V^{-\eta}\).
The comparison \(N^2/H\ge V^{1/2-o(1)}\) is a diagnostic of
the source's proof sketch at this matched scale. It is not an
application of its theorem beyond its hypotheses.

\paragraph{The exact differencing assertion and its limitation.}
With \(|x|,y\ll\sqrt Z\), Cauchy--Schwarz applied to
\[
 S=\sum_y\sum_x \mu(x^2+3y^2)\Xi(x,y)
\]
produces the two-Möbius factor
\(\mu(k+\Delta)\mu(k)\), where
\(\Delta=3(y_1^2-y_2^2)\) and \(k=x^2+3y_2^2\).
The nonzero differences have divisor-bounded multiplicities.
When the \(y\)-interval has length comparable to \(\sqrt Z\),
they have \(Z^{1-o(1)}\) distinct values in a range of order \(Z\).
This lower bound is not asserted for a shorter interval or a thin subset.
The polynomial arguments remain of size \(Z\), and their base
depends on a factorization of \(\Delta/3\). This is not the
single-Möbius sum with independent shift weights in GMPP's theorem.
A bound for this correlation is not proved here.
A possible coefficient-norm improvement mentioned below is
only an optimistic diagnostic for one Cauchy step,
not a strengthened version of the source theorem.

\section{Proof overview}
We first record the source hypotheses and distinguish the single-norm
scale \(Z\) from the product scale \(\mathcal X\asymp Z^2\).
The normalization modulo 2 converts primitive Eisenstein ideals
to the form \(x^2+3y^2\) without multiplying the Möbius argument
by four. Counting the permitted square shifts exposes the loss
in using a bound that still pays for the ambient shift interval.

The product identity and an affine-lattice count
show that its admissible small-shift region is already small.
A direct comparison of the multiplicative factor sizes then
excludes the stated Type II range. Finally we write the exact
Cauchy correlation. Although its differences cover many shifts,
the additional Möbius factor and the varying quadratic base
prevent a direct application of the independent-shift theorem.
Each conclusion is limited to the displayed transformation.

\section{Source and exact hypotheses}\label{source-and-exact-hypotheses}

Primary source:

Grimmelt, L., Merikoski, J., Pandey, M., \& Pascadi, A. (2026).
\emph{Primes and the Möbius function over quadratic polynomials}
{[}Preprint, version 1{]}. https://arxiv.org/abs/2610.41356. The
argument below uses only the indicated statements and proof steps; no
assertion about the correctness of the entire preprint is needed for the
application tests.

The source's variable named \(X\) is denoted \(V\) throughout this paper
to prevent a scale mismatch with the actual product of our inverse
norms.

\subsection{The Möbius theorem and its weighted
proof}\label{the-muxf6bius-theorem-and-its-weighted-proof}

Theorem 1.3, printed p.~3, is the first-moment assertion

\[
\sum_{h\le H}\left|\sum_{n^2+h\le V}\mu(n^2+h)\right|
\ll \nu H V^{1/2},\qquad H=V^{1/2-\nu}.
\tag{1}
\]

The proof in §12, printed pp.~57--58, explicitly begins with arbitrary
weights \(|\gamma_h|\le1\) and bounds the corresponding weighted sum.
Thus selecting a subset of shifts with a bounded weight is permitted.
The bound still contains the \textbf{ambient length \(H\)}, rather than
the cardinality of the selected subset.

The proof takes \(0<\eta<\nu/100\), and may suppose \(\nu<1/100\). Its
final paragraph obtains the stated \(o(HV^{1/2})\) near the square-root
range by fixing \(\nu_0>0\), running the argument uniformly for
\(V^{1/2-\nu_0}\le H\le V^{1/2}\), and then letting \(\nu_0\) tend to
zero. It does not supply a power-uniform estimate permitting
\(\nu=V^{-c}\) while all proof constants and thresholds remain
controlled.

\subsection{What the Type II theorem actually
estimates}\label{what-the-type-ii-theorem-actually-estimates}

Section 2.1, printed pp.~4--5, fixes a smooth function \(f\) supported
in \([1/2,1]\), with derivative bounds \(f^{(j)}\ll_j\delta^{-j}\) and
\(\delta^{-1}\le(\log V)^{O(1)}\). It takes

\[
R_{\mathrm G}=V^{\eta/100}.
\]

Section 2.2, Definition 1, printed p.~5, defines a smooth square
indicator \(\sigma_1(k)=f(k/V)\mathbf1_{k=\square}\), a major-arc model
\(\sigma_0(k)\), their difference \(\sigma^\flat\), and

\[
w_t^{(\gamma)}=\sum_{k+h=t}\gamma_h\sigma^\flat(k),
\qquad \gamma_h=0\ (h>H),\quad |\gamma_h|\le V^{o(1)}.
\tag{2}
\]

Proposition 2.4, printed p.~6, assumes

\[
H=V^{1/2-\nu}\le V^{1/2+o(1)},\quad
\nu\le\tfrac16-2\eta,\quad
V^{\nu+\eta}<N<V^{1/4-\nu/2-\eta},\quad MN=V.
\tag{3}
\]

For divisor-bounded \(\alpha_m,\beta_n\), with \(\beta_n\) supported on
squarefree integers, it gives

\[
\sum_{m\asymp M}\sum_n\alpha_m\beta_n w_{mn}^{(\gamma)}
\ll V^{o(1)}HV^{1/2}R_{\mathrm G}^{-1/4}
=V^{o(1)}HV^{1/2-\eta/400}.
\tag{4}
\]

This is an estimate for the difference from the major-arc model, with
separable coefficients \(\alpha_m\beta_n\gamma_h\). It is not an
estimate for an arbitrary bounded weight in \((m,n,h)\), and it is not
an assertion for every \(M,N\) with \(MN=V\).

The exact locators used below are §2.5, printed pp.~8--9, for the
diagnostic range calculation; §8.1, pp.~36--37, for the first Cauchy
step and diagonal; §8.2, pp.~38--39, for pseudo-diagonals; §8.3,
pp.~39--40, for removal of the shift weights; and §12, equations
(12.1)--(12.4), pp.~57--58, for the final Möbius theorem.

\section{The actual ideal sum and a common scale
dictionary}\label{the-actual-ideal-sum-and-a-common-scale-dictionary}

Let \(K=\mathbf Q(\omega)\), with the fixed cubic character and deletion
of 3 and 269 specified above. The elementary content decomposition
writes each inverse ideal as

\[
\mathfrak d_i=(r_i)\mathfrak b_i,
\qquad \mathfrak b_i\text{ primitive},
\qquad m(\mathfrak d_i)=\mu(r_i)^2\chi_{-3}(r_i)
m(\mathfrak b_i)\mathbf1_{(r_i,N\mathfrak b_i)=1}.
\]

The retained region still includes \(r_1=r_2=1\). Fixing the short
direct ideal and considering that primitive corner gives a natural
difficult dyadic block

\[
\mathcal B=\sum_{N\mathfrak b_1\asymp Z}
\sum_{N\mathfrak b_2\asymp Z}
\mu_K(\mathfrak b_1)\psi(\mathfrak b_1)
\mu_K(\mathfrak b_2)\psi(\mathfrak b_2)
F\!\left(\frac{N\mathfrak b_1N\mathfrak b_2}{\mathcal X}\right),
\qquad \mathcal X\asymp Z^2.
\tag{5}
\]

Fixed smooth caps and fixed residue restrictions may be included.
Polylogarithmic changes of the dyadic endpoints do not alter the power
calculations. Other retained content blocks can have unbalanced
primitive norms; the claim here concerns a primitive balanced corner
that the content deletion does not remove. Bounding the whole sum need
not logically require bounding each block separately, but a proof that
treats these blocks separately must address this one.

There is an exact convenient coordinate normalization. A primitive ideal
with nonzero ideal Möbius coefficient has odd squarefree norm: the prime
2 is inert and cannot divide a primitive ideal. Modulo 2, the Eisenstein
ring is \(\mathbf F_4\). Among the six unit multiples of a generator,
exactly two, differing by sign, reduce to 1. Such a generator has the
form

\[
a+b\omega=x+y\sqrt{-3},\qquad b=2y,\quad x=a-b/2,
\]

with \(a\) odd. Choosing \(y>0\) removes the remaining sign ambiguity
except for the unit ideal. Its norm is

\[
Q(a,b)=a^2-ab+b^2
=\frac{(2a-b)^2+3b^2}{4}=x^2+3y^2.
\tag{6}
\]

For a nonzero Möbius term, the conditions \((x,y)=1\) and \(x+y\) odd
hold automatically: a common divisor makes its square divide the norm,
and two odd coordinates give a norm divisible by 4. Thus there is no
illegitimate replacement of \(\mu(Q)\) by \(\mu(4Q)\). The latter would
be identically zero and must not be used.

The finite character is a function of \(x,y\) modulo a fixed modulus;
splitting into finitely many residue classes isolates it. GMPP does not
explicitly state the needed theorem for every such restriction on the
square-root coordinate. For the negative applicability tests below, we
grant the most favorable fixed-modulus and smooth-weight extensions: the
density and scale failures persist even with those extra extensions. No
such unprinted extension is being imported as a theorem.

The scale dictionary is:

\begin{longtable}[]{@{}
  >{\raggedright\arraybackslash}p{(\columnwidth - 6\tabcolsep) * \real{0.2000}}
  >{\raggedleft\arraybackslash}p{(\columnwidth - 6\tabcolsep) * \real{0.2667}}
  >{\raggedleft\arraybackslash}p{(\columnwidth - 6\tabcolsep) * \real{0.2667}}
  >{\raggedleft\arraybackslash}p{(\columnwidth - 6\tabcolsep) * \real{0.2667}}@{}}
\toprule\noalign{}
\begin{minipage}[b]{\linewidth}\raggedright
Attempt
\end{minipage} & \begin{minipage}[b]{\linewidth}\raggedleft
GMPP polynomial scale \(V\)
\end{minipage} & \begin{minipage}[b]{\linewidth}\raggedleft
Typical actual shift
\end{minipage} & \begin{minipage}[b]{\linewidth}\raggedleft
GMPP allowed \(H\)
\end{minipage} \\
\midrule\noalign{}
\endhead
\bottomrule\noalign{}
\endlastfoot
One norm \(x^2+3y^2\asymp Z\) & \(Z\) & \(3y^2\asymp Z\) & at most
\(Z^{1/2+o(1)}\) \\
Product of two norms & \(\mathcal X\asymp Z^2\) &
\(3(x_1y_2-x_2y_1)^2\asymp Z^2\) & at most \(Z^{1+o(1)}\) \\
Difference \(3(y_1^2-y_2^2)\) & arguments of each Möbius factor remain
size \(Z\) & difference ranges through \(O(Z)\) & at scale \(Z\), at
most \(Z^{1/2+o(1)}\) \\
\end{longtable}

Calling a difference of size \(Z\) an admissible shift at product scale
\(Z^2\) does not change the sizes of the Möbius arguments. That mismatch
is one of the failures below.

\section{Selecting square shifts: the exact density
loss}\label{selecting-square-shifts-the-exact-density-loss}

First work at individual norm scale \(V=Z\), and restrict
\(3y^2\le H=V^{1/2-\nu}\). There are

\[
B\asymp H^{1/2}=V^{1/4-\nu/2}
\]

possible positive \(y\), and \(O(V^{1/2})\) square-root coordinates
\(x\). The selected-slice trivial bound is therefore

\[
B V^{1/2}\asymp V^{3/4-\nu/2}.
\tag{7}
\]

Taking \(\gamma_h\) supported on \(h=3y^2\) is allowed in the weighted
proof of (1), but gives only

\[
\nu HV^{1/2}=\nu V^{1-\nu}.
\tag{8}
\]

Relative to (7), the loss is

\[
\nu H^{1/2}=\nu V^{1/4-\nu/2}.
\tag{9}
\]

For fixed \(\nu>0\), this grows polynomially. The limiting \(o(1)\) in
Theorem 1.3 has no demonstrated polynomial rate capable of defeating
that loss.

Even isolating just the Type II residual (4), the sparse shift bound
loses

\[
H^{1/2}V^{-\eta/400}
=V^{1/4-\nu/2-\eta/400+o(1)}
\tag{10}
\]

against the selected support. In the stated range
\(0\le\nu\le1/6-2\eta\), one has \(\eta\le1/12\). Thus
\(\eta/400\le1/4800\), whereas \(1/4-\nu/2\ge1/6\). The loss cannot be
overcome by the saving supplied by this proposition. In the actual §12
proof, \(\eta<\nu/100\) and \(\nu<1/100\) make the available saving
smaller still.

\subsection{A black-box nonimplication, including two
copies}\label{a-black-box-nonimplication-including-two-copies}

This is not merely a weakness of a particular triangle inequality.
Suppose an abstract array \(S_h\) obeys \(|S_h|\le A\) and

\[
\sum_{h\le H}|S_h|\le EHA.
\]

Let \(\mathcal S\) contain \(B\) square shifts. If \(B\le EH\), the
array

\[
S_h=A\mathbf1_{h\in\mathcal S}
\]

obeys every such bounded-weight first-moment estimate but has no
cancellation on the selected set. For (1), \(E\asymp\nu\), and for the
optimistic use of (4), \(E=V^{-\eta/400+o(1)}\). The inequality
\(B\le EH\) holds in both ranges above for large \(V\).

Taking two copies gives a tensor array whose sum on
\(\mathcal S\times\mathcal S\) is exactly \(A^2B^2\). Replacing the pair
by the difference of its indices can make the difference set dense, but
its convolution remains nonnegative and has the same total mass.
Consequently the stated first-moment bound, tensoring, and
Cauchy--Schwarz alone cannot force cancellation on the paired square
supports. This is a logical nonimplication for the supplied estimates,
not a model for the actual values of the Möbius function.

\subsection{Retaining the shift norm in the second Cauchy
step}\label{retaining-the-shift-norm-in-the-second-cauchy-step}

The actual second Cauchy step is visible in §8.3, printed p.~39. It
removes \(\gamma_{h_1}\overline\gamma_{h_2}\) by bounding the
corresponding two-dimensional coefficient norm and then extends a
nonnegative sum to a smooth full lattice. In the favorable coprime part,
retaining a support of size \(B\), rather than replacing it by \(H\),
improves the fourth-root final estimate by at most

\[
(B/H)^{1/2}.
\]

Thus an optimistic version of the relevant off-diagonal bound would be

\[
V^{1/2}(BH)^{1/2}V^{-\eta/400+o(1)}.
\tag{11}
\]

Against \(BV^{1/2}\), this still loses

\[
H^{1/4}V^{-\eta/400}
=V^{1/8-\nu/4-\eta/400+o(1)}.
\tag{12}
\]

Here \(1/8-\nu/4\ge1/12\), far larger than the available \(\eta/400\).
Formula (11) is a deliberately favorable diagnostic for that Cauchy
step, \textbf{not} a claimed strengthened Proposition 2.4: the source
handles the pseudo-diagonals and other terms separately, and they have
not been proved to enjoy this refinement. Even this favorable gain does
not close the sparse-support loss.

\section{What portion of our balanced sum lies in this one-norm
range?}\label{what-portion-of-our-balanced-sum-lies-in-this-one-norm-range}

At norm scale \(Z\), the ordinary two-dimensional shell has \(O(Z)\)
points. Restricting \(3y^2\le Z^{1/2-\nu}\) leaves only

\[
O(Z^{3/4-\nu/2})
\]

points, by counting \(O(\sqrt Z)\) choices for \(x\) and
\(O(Z^{1/4-\nu/2})\) for \(y\). In the pair block (5), leaving the other
ideal unrestricted therefore gives the absolute bound

\[
O(Z^{7/4-\nu/2})
=O(\mathcal X^{7/8-\nu/4}).
\tag{13}
\]

If both ideals are in these small-coordinate slices, the absolute bound
is

\[
O(Z^{3/2-\nu})=O(\mathcal X^{3/4-\nu/2}).
\tag{14}
\]

These elementary estimates are for the indicated pieces only. The
complementary part, where both relevant coordinates have their natural
square-root size, remains. Applying GMPP in its original polynomial
scale therefore addresses a region already at or below the seven-eighths
scale by counting and does not address the bulk of the balanced
expression.

\section{Multiplying the two ideals: a square shift
remains}\label{multiplying-the-two-ideals-a-square-shift-remains}

For the coordinate representatives \(z_i=x_i+y_i\sqrt{-3}\), complex
multiplication and conjugation give the exact identity

\[
(x_1^2+3y_1^2)(x_2^2+3y_2^2)
=(x_1x_2+3y_1y_2)^2
 +3(x_1y_2-x_2y_1)^2.
\tag{15}
\]

The alternate sign version is obtained by conjugating one factor. The
first square-root coordinate has size \(Z\), and the determinant

\[
w=x_1y_2-x_2y_1
\]

also has typical size \(Z\). Hence the polynomial scale is now correctly
\(V=\mathcal X\asymp Z^2\), but the shift is \(h=3w^2\asymp V\), not
\(h\le V^{1/2}\). The product construction retains a square-supported
shift; it does not create a full interval of shifts.

\subsection{Exact counting in the part with an admissibly small
shift}\label{exact-counting-in-the-part-with-an-admissibly-small-shift}

Suppose \(3w^2\le H\), with \(H\le V^{1/2+o(1)}\). For each fixed
primitive \((x_1,y_1)\) with norm comparable to \(Z\), the equation

\[
x_1y_2-x_2y_1=w
\]

has solutions in one affine lattice of step \((x_1,y_1)\), because
\((x_1,y_1)=1\). That step has Euclidean length comparable to
\(\sqrt Z\), and the norm-\(Z\) shell for the second vector has diameter
\(O(\sqrt Z)\). Thus there are \(O(1)\) second vectors for each integer
\(w\). There are \(O(Z)\) first vectors. The number of pairs with
\(|w|\le\sqrt{H/3}\) is therefore

\[
O\bigl(Z(1+\sqrt H)\bigr).
\tag{16}
\]

For \(H=V^{1/2-\nu}\), this is

\[
O\bigl(V^{1/2}+V^{3/4-\nu/2}\bigr).
\tag{17}
\]

The bounded character and Möbius factors cannot increase that estimate.
A finite choice of unit rotations or conjugation changes it only by a
constant. The portion of the product representation in the GMPP shift
range is therefore an already small near-parallel region. Most of the
pairs have \(|w|\asymp Z\) and lie outside that range.

\subsection{Coefficient matching is also
nontrivial}\label{coefficient-matching-is-also-nontrivial}

The expression (5) has the smooth product weight
\(F(q_1q_2/\mathcal X)\), not a freely inserted square-indicator model
\(w^{(\gamma)}_{q_1q_2}\). Identity (15) gives representations of a
product norm, but their multiplicities, balanced factorization
restrictions, and cubic phases are tied to its two factors. They do not
automatically factor as the independent \(\gamma_h\) in (2).

Adding an artificial dense average of shifts does not solve this: its
major-arc model has a principal term of size \(H/(2\sqrt{q_1q_2})\), so
at \(H\asymp\sqrt{\mathcal X}\) it reproduces a smooth multiple of the
original bilinear expression. Bounding the model then requires the same
inverse-Hecke cancellation. In §12 of GMPP the model is handled for
ordinary Möbius coefficients by different arithmetic; that step is not
an estimate for the present fixed cubic inverse coefficients.

\section{The Type II range, with the polynomial scale matched
correctly}\label{the-type-ii-range-with-the-polynomial-scale-matched-correctly}

Even granting, favorably and without claiming it, that a product-norm
block could be represented with all the separability and model
requirements of (2), its multiplicative factors would be

\[
V=\mathcal X\asymp Z^2,\qquad M\asymp N\asymp Z\asymp V^{1/2}.
\tag{18}
\]

Proposition 2.4 instead requires

\[
N<H^{1/2}V^{-\eta}\le V^{1/4+o(1)}.
\tag{19}
\]

Thus the actual balanced factor exceeds the available upper range by a
factor \(V^{1/4-o(1)}\). Interchanging the factors changes nothing.
Fixed powers of logarithms in the caps cannot bridge this power gap.
Small rational contents do not repair this for the primitive corner
\(r_1=r_2=1\).

The source's §2.5 sketch makes the origin of the upper endpoint
explicit. After its two Cauchy steps it compares a reference
fourth-moment size \(H^2N^2\) with a large-sieve contribution \(HN^4\);
their ratio is

\[
N^2/H.
\tag{20}
\]

With the correctly matched values (18) and \(H\le V^{1/2+o(1)}\), this
ratio is at least \(V^{1/2-o(1)}\), rather than a negative power. Taking
the fourth root would lose \(V^{1/8-o(1)}\) relative to that comparison.
This is a diagnostic of the \textbf{published sketch's range}, not an
established upper bound for our sum and not a license to apply its
theorem outside (3). The rigorous nonapplicability is already the
incompatible inequality (19).

\section{Differencing can make the shifts dense, but changes the
arithmetic
object}\label{differencing-can-make-the-shifts-dense-but-changes-the-arithmetic-object}

Take a rectangular subregion with \(|x|\ll\sqrt Z\),
\(1\le y\ll\sqrt Z\), and first suppress the harmless fixed residue
weights and smooth caps. Write

\[
S=\sum_y\sum_x f(x,y),\qquad
f(x,y)=\mu(x^2+3y^2)\,\Xi(x,y),
\]

where \(\Xi\) represents the bounded finite character and smooth
restrictions. Cauchy in the \(x\)-coordinate gives the exact expansion

\[
|S|^2\le O(\sqrt Z)
\sum_{y_1,y_2}\sum_x
\mu(x^2+3y_1^2)\mu(x^2+3y_2^2)
\Xi(x,y_1)\overline{\Xi(x,y_2)}.
\tag{21}
\]

The complete right side is real and nonnegative; its individual terms
need not be. Set

\[
\Delta=3(y_1^2-y_2^2),\qquad k=x^2+3y_2^2.
\]

The arithmetic factor becomes

\[
\mu(k+\Delta)\mu(k),
\quad k\text{ constrained to }x^2+3y_2^2.
\tag{22}
\]

For nonzero \(\Delta\), the number of positive pairs \((y_1,y_2)\) with
that difference is bounded by a divisor function, since

\[
(y_1-y_2)(y_1+y_2)=\Delta/3.
\]

For this lower bound, assume the positive \(y\)-interval has length
comparable to \(\sqrt Z\). There are then order \(Z\) off-diagonal
pairs in its square, and each
difference has at most \(Z^{o(1)}\) preimages. Thus the difference set
has \(Z^{1-o(1)}\) distinct values in a range of order \(Z\). The second
coordinate really does generate a dense set on the level of powers. This
is the strongest part of the proposed analogy. No such lower bound
is asserted for an arbitrary short interval or thin subset.

However, (22) is a \textbf{two-Möbius correlation on a quadratic
subsequence}, with the base quadratic shift \(3y_2^2\) tied to a
factorization of \(\Delta/3\). It is not \(\mu(x^2+h)\) with an
independent coefficient \(\gamma_h\). Fixing \(y_2\) makes the allowed
\(\Delta\) a translate of the sparse square set again. Summing \(y_2\)
makes the additional factor \(\mu(x^2+3y_2^2)\) depend jointly on the
square-root variable and the shift data. The theorem's separable bounded
shift weight does not allow that dependence.

This distinction is essential, not a harmless unproved uniformity. If
the theorem allowed an arbitrary bounded weight \(\gamma_{h,x}\),
choosing it equal to \(\mu(x^2+h)\) would turn the Möbius sum into a
nonnegative squarefree count. Thus arbitrary dependence on the
square-root variable cannot be added to (1). The actual extra factor in
(22) has arithmetic structure, but exploiting that structure is a new
correlation problem.

Here the contradiction can be made quantitative without assuming
a theorem on squarefree values of a fixed quadratic polynomial.
Take \(H=\lfloor\sqrt V\rfloor\) and
\(n\in[a\sqrt V,b\sqrt V]\), for fixed \(0<a<b<1\).
For large \(V\), all the values \(n^2+h\), \(1\le h\le H\), lie below \(V\).
For each \(n\), the number divisible by some prime square is at most
\[
 H\sum_p p^{-2}+\pi(\sqrt V).
\]
Indeed, \(p^2\mid n^2+h\) selects at most \(H/p^2+1\) values of \(h\),
and only \(p\le\sqrt V\) can occur.
Now \(\sum_p p^{-2}\le\sum_{m\ge2}m^{-2}<1\), while
\(\pi(t)\ll t/\log t\).
The latter elementary estimate follows by observing that every prime
in \((m,2m]\) divides \(\binom{2m}{m}\le4^m\), and then summing
the resulting bounds over dyadic intervals.
Thus a fixed positive fraction of the \(h\)'s survives for each \(n\).
It follows that
\(\sum_{h,n}\mu(n^2+h)^2\gg HV^{1/2}\), contrary to a uniform
\(o(HV^{1/2})\) conclusion with arbitrary bounded
\(\gamma_{h,n}\). A fixed nonnegative smooth weight bounded below
on the chosen \(n\)-interval gives the same conclusion.

Applying Cauchy directly to the already paired sum creates four Möbius
factors rather than eliminating this issue. Reindexing by
\(3(y_1^2-y_2^2)\) before Cauchy leaves a varying base shift and the
second Möbius factor; it does not turn the two factors into one legal
instance of (1).

\subsection{Why the product scale cannot be borrowed for
(22)}\label{why-the-product-scale-cannot-be-borrowed-for-22}

The difference \(|\Delta|\ll Z\) happens to be of size at most
\(\mathcal X^{1/2}\), but each argument in (22) is still of size \(Z\),
not \(\mathcal X\asymp Z^2\). At their true polynomial scale \(V=Z\),
the difference range extends to \(V\), beyond (1)'s range.

If one artificially uses \(V=Z^2\), the square-root coordinate in (22)
has length \(Z^{1/2}=V^{1/4}\), instead of \(V^{1/2}\). In the smooth
formulation of §2.1 the argument \(x^2/V\) lies near \(1/Z\), outside
the fixed support \([1/2,1]\). A cutoff that isolates this region after
that normalization needs derivatives of size \(Z^j\), whereas the
theorem only permits polylogarithmic derivative scales. The unsmoothed
summatory form likewise pays the larger factor \(V^{1/2}=Z\), losing
\(Z^{1/2}\) against the actual row length. Returning to the true scale
removes that loss but restores the excessive shift range. Neither choice
yields a valid useful import.

\subsection{The diagonal budget and what a new mean would
need}\label{the-diagonal-budget-and-what-a-genuinely-new-mean-would-need}

The diagonal \(y_1=y_2\) in (21) is at most \(O(Z)\) before the outer
Cauchy factor, and therefore at most \(O(Z^{3/2})\) after it. If a new
argument controlled the full off-diagonal at that same level, it would
give \(S\ll Z^{3/4}\); paired blocks of such sums could reach a
three-quarter exponent in the product scale. This is a conditional
calculation of a useful target, not a proved estimate. The
GMPP single-Möbius average does not bound that off-diagonal.
Cancellation between the diagonal and off-diagonal is also not ruled
out, so this budget is not a lower-bound obstruction to stronger
methods.

\section{A second-ideal Cauchy step without polynomial
reparametrization}\label{a-second-ideal-cauchy-step-without-polynomial-reparametrization}

Grouping the primitive coefficients by norm gives a scalar sequence

\[
h(q)=\sum_{N\mathfrak b=q,\,\mathfrak b\ {
m primitive}}
\mu_K(\mathfrak b)\psi(\mathfrak b),\qquad |h(q)|\le d(q),
\]

supported on squarefree split norms. A dyadic block is

\[
\sum_{m\asymp Z}\sum_{n\asymp Z}h(m)h(n)F(mn/\mathcal X).
\tag{23}
\]

Cauchy in one factor produces correlations in the other with a smooth
kernel depending on ratios of norms. For instance, before grouping ideal
multiplicities, its inner kernel has the form

\[
\sum_{\mathfrak b_1}
F(N\mathfrak b_1N\mathfrak b_2/\mathcal X)
\overline{F(N\mathfrak b_1N\mathfrak b'_2/\mathcal X)}.
\tag{24}
\]

For a nonnegative shell weight this is nonnegative and has many large
off-diagonal entries. There is no new free additive shift or oscillating
spectral parameter in (24). Mellin separation expresses (23) as products
of length-\(Z\) inverse polynomials integrated over a short frequency
band. Generic mean-square estimates on that band retain their length
term \(Z\); the second copy does not turn the band into a long
independent average.

This is the same structural issue visible in the fixed finite-field
cubic kernel, which has rank one under multiplication. It is an
obstruction to these direct norm-only Cauchy and generic large-sieve
steps, not to every possible use of the arithmetic coefficients.

\section{Outcome and exact outstanding
input}\label{outcome-and-exact-outstanding-input}

The investigation establishes the following application failures with
their losses:

\begin{enumerate}
\def\labelenumi{\arabic{enumi}.}
\tightlist
\item
  \textbf{Selecting square shifts:} the permitted bounded-weight theorem
  retains an ambient-density loss \(H^{1/2}\); its Type II power saving
  is far too small. Even a favorable coefficient-norm refinement at the
  second Cauchy step leaves loss \(H^{1/4}\).
\item
  \textbf{Using the individual norm scale:} the theorem's small-shift
  region occupies a piece of the paired sum already bounded by
  \(\mathcal X^{7/8+o(1)}\) through absolute counting.
\item
  \textbf{Multiplying the ideals:} the new shift remains a square,
  usually of product-norm size. Its admissible small-shift part has the
  stronger elementary bound \(\mathcal X^{3/4+o(1)}\).
\item
  \textbf{Importing the stated Type II proposition at product scale:}
  balanced factors are a power \(V^{1/4-o(1)}\) beyond its upper
  endpoint, even after favorable coefficient-matching assumptions. The
  \(N^2/H\) comparison is explicitly tied to this correctly matched
  scale and only to the source's sketch.
\item
  \textbf{Differencing square shifts:} it does generate many distinct
  shifts, but requires a uniform two-Möbius correlation with a
  shift-dependent quadratic base. That is a different analytic input
  from the cited theorem.
\end{enumerate}

No new bound for the full retained cubic sum has been proved. A route
through these transformations would need a new estimate for the precise
correlation in (22), or a sparse-shift/variable-form dispersion result
strong enough to retain the arithmetic support and valid in the balanced
factor range. The GMPP theorem as stated does not supply either. Whether
its proof can be substantially redesigned to do so remains open in this
work.


\section*{Method and availability}
The mathematics and text were produced by GPT-6 Astra in ChatGPT Pro mode.
The cited GMPP theorem is an external input. The sparse-support
nonimplications, ideal normalization, lattice counts, scale comparisons,
and exact differencing identities are proved here under their stated
hypotheses. The conclusions concern the specified applications only.
The accompanying records distinguish a targeted whole-record check from
a separate statement-only reconstruction of the elementary claims by
the same model. No human review or formal verification is claimed.
Sources and verification records are available at
\href{\workbenchurl}{the workbench repository}, with version records at
\href{\recordurl}{the research record}.
The original text and proof organization are dedicated to the public
domain under CC0 1.0; the cited preprint retains its own terms.

\begin{thebibliography}{9}
\bibitem[Grimmelt et al.(2026)]{GMPP}
Grimmelt, L., Merikoski, J., Pandey, M., \& Pascadi, A. (2026).
\emph{Primes and the Möbius function over quadratic polynomials}
[Preprint, version 1]. arXiv.
\url{https://arxiv.org/abs/2610.41356}.
Theorem 1.3, p.~3; Definition 1, p.~5; Proposition 2.4, p.~6;
\(\S2.5\), pp.~8--9; \(\S\S8.1\)--8.3, pp.~36--40;
\(\S12\), pp.~57--58.
\end{thebibliography}
\end{document}
